\subsection{问题一：半平面交会与定位区域直径}
\label{sub:model_p1}
\subsubsection{示向约束及区域状态判定}
设第$i$个检测点为$S_i$，示向度为$\theta_i$，误差界为$\varepsilon=1^\circ$，定义边界方向
\begin{equation}
 u_i^{\pm}=\bigl(\cos(\theta_i\pm\varepsilon),\sin(\theta_i\pm\varepsilon)\bigr).
\end{equation}
记二维叉积为$[a,b]=a_xb_y-a_yb_x$，则该次观测对应的角扇区为
\begin{equation}
 W_i=\{X:[u_i^-,X-S_i]\ge0,\ [u_i^+,X-S_i]\le0\},\qquad
 P=\bigcap_{i=1}^{n}W_i=\{X:AX\le b\}.
 \label{eq:angular_intersection}
\end{equation}
向量表示可直接处理跨越$0^\circ$的示向度。这里计算题目所指的纯角度交会区域；目标圆、接收距离等附加信息在第二问中另行使用。

为避免将空集或无界区域误判为有限多边形，首先枚举不平行边界线的交点，保留满足全部约束者；若未找到可行交点，再检查原点及原点在各边界上的正交投影。该判定是完备的：非空闭凸集距原点最近的点，必为原点、单条有效边界上的垂足或两条独立边界的交点。其次，对非空交集检验
\begin{equation}
 \exists d\ne0,\quad Ad\le0.
 \label{eq:recession}
\end{equation}
若成立，则从任一可行点沿$d$无限延伸仍保持可行，区域无界；否则区域有界。二维实现仅需检查各非零约束法向量对应的两条切向方向。算法不引入人为有限边框。对有界交集取凸包，根据顶点数区分单点、线段和多边形。

\subsubsection{直径计算与覆盖性判据}
设有界区域的顶点为$v_1,\ldots,v_k$。任意$p,q\in P$均可写成顶点的凸组合，从而
\begin{equation}
 \|p-q\|\le\sum_{i,j}\alpha_i\beta_j\|v_i-v_j\|
 \le\max_{i,j}\|v_i-v_j\|,
 \qquad D(P)=\max_{i,j}\|v_i-v_j\|.
 \label{eq:vertex_diameter}
\end{equation}
因此只需求最远顶点对。本文采用旋转卡壳方法\cite{toussaint}，在逆时针凸包上依次更新对踵支撑点；遇到平行支撑边时同时检查两端点。直径阶段需要$O(k)$次算术操作，边界交点枚举及约束检验仍为$O(m^3)$，其中$m$为约束数，故不能将完整算法视为线性复杂度。程序对已构造的半平面系数使用有理数进行几何判断；三角函数系数本身仍存在机器舍入误差。

若$a,b$为最远点对，任何半径$D/2$且覆盖$a,b$的圆，其圆心只能是$c_D=(a+b)/2$。由圆盘的凸性，存在直径为$D$的覆盖圆当且仅当
\begin{equation}
 \max_j\|v_j-c_D\|\le D/2.
 \label{eq:diameter_circle}
\end{equation}
该条件并非总成立。构造等边三角形$A=(0,0)$、$B=(40,0)$、$C=(20,20\sqrt3)$，取三个检测点
\begin{equation}
 S_1=(-1000,0),\quad S_2=(540,-500\sqrt3),\quad
 S_3=(520,520\sqrt3),
\end{equation}
示向度依次为$1^\circ,121^\circ,241^\circ$。三条下边界对应三角形的有向边，而另一侧最大偏角为$\arctan(20\sqrt3/1020)<2^\circ$，故三扇区交集恰为该三角形。真实源可取其重心，接收半径取1200米，与题设相容。此时$D=40$米，而最小覆盖圆半径为$40/\sqrt3=23.0940$米，故直径为40米的圆不能覆盖定位区域。

最小覆盖圆由至多三个边界点确定\cite{welzl}。本文枚举两点直径圆与三点外接圆，选择覆盖全部顶点的最小者；该确定性实现最坏需要$O(k^4)$次算术操作。记其半径为$R_*$，平面上的Jung不等式\cite{jung}给出
\begin{equation}
\frac{D}{2}\le R_*\le\frac{D}{\sqrt3}.
 \label{eq:jung}
\end{equation}
下界由最远点对立即得到；上界亦可由支撑点说明：两点支撑时$R_*=D/2$，三点支撑时圆心在支撑三角形内，其最大角$\gamma\in[60^\circ,90^\circ]$，由正弦定理$R_*\le D/(2\sin\gamma)\le D/\sqrt3$。
因此，一次20米光学定位能够保证覆盖全部候选源的准确判据是$R_*\le20$；仅知直径时，$D\le20\sqrt3$是充分条件，而$D\le40$仅为必要条件。

\begin{table}[htbp]
 \centering
 \caption{第一问自建算例的几何结果}
 \label{tab:p1_geometry}
 \begin{tabular}{lrrc}
 \hline
 算例 & $D$/米 & $R_*$/米 & 直径为$D$的圆可覆盖\\
 \hline
 双站正交交会 & 4.9385 & 2.4693 & 是\\
 边长20米的三角形 & 20.0000 & 11.5470 & 否\\
 边长40米的三角形 & 40.0000 & 23.0940 & 否\\
 较宽基线交会 & 31.4797 & 15.7398 & 是\\
 \hline
 \end{tabular}
\end{table}
平行发散观测与不相容观测分别得到无界区域和空集。以上均为自行构造的数学算例，原题未提供第一问的具体观测数据。

\subsection{问题二：保证收信的第二检测点选择}
\label{sub:model_p2}
\subsubsection{首验信息与最坏直径目标}
题设仅给出测向误差界，本文以与观测及题设约束相容的全部源位置表示定位结果，不额外指定其概率分布。
取首次检测点$S_1$，沿首示向方向建立局部正交坐标，令$u=(\cos\theta,\sin\theta)$、$n=(-\sin\theta,\cos\theta)$，第二点写为$S(q)=S_1+a u+b n$，其中$q=(a,b)$。首次正常读数对应的候选源集合为
\begin{equation}
 U_1=\{G:\|G\|\le1800,\ 5<\|G-S_1\|\le1500,
 \ |\operatorname{wrap}(\arg(G-S_1)-\theta)|\le\varepsilon\}.
\end{equation}
这里$\operatorname{wrap}$表示取环形角差。固定首验后，记$\Psi(q)$为全部相容的第二次正常读数，收到读数$\psi$后保留下来的源位置集合为
\begin{equation}
 \begin{aligned}
 U(q,\psi)=U_1\cap\{G:\;&5<\|G-S(q)\|\le1500,\\
 &|\operatorname{wrap}(\arg(G-S(q))-\psi)|\le\varepsilon\}.
 \end{aligned}
\end{equation}
两次正常读数与同一未知半径相容，当且仅当$\max\{1000,\|G-S_1\|,\|G-S(q)\|\}\le1500$，故上述距离条件已包含这一信息。采用Tokekar和Isler的最坏扇区交集直径准则\cite{tokekar}，并将其条件化于已知首验，建立
\begin{equation}
 J(q)=\sup_{\psi\in\Psi(q)}\operatorname{diam}U(q,\psi),\qquad
 \min_{q\in\mathcal F}J(q).
 \label{eq:p2_objective}
\end{equation}
该指标与第一问一致，衡量最不利合法读数下的剩余位置歧义，无需假设源位置先验或误差概率分布。
直径刻画候选位置间的最大分离，可反映面积较小但形状细长的区域仍具有的位置歧义。候选点须对全部相容源位置及接收半径保证收信，这构成鲁棒可行性约束\cite{bertsimas2011}；在该约束下再比较最坏定位损失。以下分别推导安全候选域及两类评分上界。

\subsubsection{四圆盘安全候选域}
设第一距离为$r$、未知接收半径为$R\in[1000,1500]$。首次收信还意味着$R\ge r$，故保证第二点对所有相容半径仍能收信的条件为
\begin{equation}
 \|S(q)-G\|\le\max(1000,r).
\end{equation}
令$u_\pm=(\cos\varepsilon,\pm\sin\varepsilon)$，暂将首验放松为完整环扇形$5<r\le1500$、$|\phi|\le\varepsilon$，则连续安全区域可写为
\begin{equation}
 \mathcal C_{\rm safe}
 =\bigcap_{d\in\{5,1000\},\ s\in\{-,+\}}
 \{q:\|q-du_s\|\le1000\}.
 \label{eq:four_disks}
\end{equation}
证明如下：$d=1000$的两个圆盘给出$q\cdot u_\pm\ge\|q\|^2/2000\ge0$，因此$q\cdot u_\phi$在允许角度内的最小值位于角端点。对$5\le r\le1000$，距离平方关于$r$为凸函数，只需检验两个径向端点；对$1000\le r\le1500$，要求化为$\|q\|^2-2r q\cdot u_\phi\le0$，最严格处为$r=1000$。由此得到四个圆盘。目标圆进一步限制源的位置，故该区域对实际问题仍安全，但不一定最大。

若直接横向移动，取真实源$(1000,0)$、$R=1000$，任意第二点$(0,b)$、$b\ne0$均会失收，说明改善交会夹角必须服从收信约束。为统一比较第二次正常测向，本文采用
\begin{equation}
 \mathcal F=\mathcal C_{\rm safe}\cap
 \{q:\|q\|\le L,\ |b|\cos\varepsilon-a\sin\varepsilon>5\}.
 \label{eq:p2_feasible}
\end{equation}
最后一项保守排除了第二点进入5米近区的情况；实际遇到近区反馈可直接清除。默认$L=1100$米，四圆盘本身已蕴含$\|q\|\le1005$米。1800米目标圆仅约束干扰源，第二检测点可以位于圆外。

\subsubsection{评分方法一：解析条带上界}
参照Tokekar和Isler对测向扇区交集作几何外包、进而界定最坏定位直径的思路\cite{tokekar}，本文结合距离上界，将两次测向扇区分别放松为条带。由条带半宽及交会夹角的下界，推导解析评分上界$A(q)$。记
\begin{align}
 w_1&=1500\sin\varepsilon,\quad h_0=|b|-w_1,\quad
 k_0=\max\{|a-5\cos\varepsilon|,|a-1500|\},\notag\\
 \beta_0&=\arctan2(h_0,k_0)-\varepsilon,\quad
 r_{2,\max}=\max_{r\in\{5,1500\},\ s\in\{-,+\}}\|q-r u_s\|,
 \quad w_2=r_{2,\max}\sin\varepsilon.
\end{align}
当$h_0>0$且$\beta_0>0$时，两条示向中心直线的锐角不小于$\beta_0$，定位区域包含于半宽分别为$w_1,w_2$的条带交集，因而
\begin{equation}
 J(q)\le A(q)=
 \frac{2\sqrt{w_1^2+w_2^2+2w_1w_2\cos\beta_0}}{\sin\beta_0}.
 \label{eq:analytic_bound}
\end{equation}
否则置$A(q)=+\infty$，表示该解析放松未给出有限保证。

\subsubsection{评分方法二：几何外包上界}
沿用Tokekar和Isler以相容测向扇区交集的最坏直径评价定位质量的方法\cite{tokekar}，本文进一步保留目标圆和距离约束，对第二读数分格并同步增宽约束，构造覆盖连续合法读数的几何外包；取各单元外包直径的最大值得到评分上界$B(q)$。

具体地，将首扇区、目标圆、首距离上界及有效投影下界求交，再加入第二距离上界圆，形成与第二读数无关的外包$Q(q)$。圆以$M$个外切半平面放松，$M=360$时1800米圆的最大径向膨胀约0.06854米；该值并非全部交会顶点的位置误差上限。两次5米排除圆以有效前向投影下界放松，保持外包为凸多边形。第二读数后的区域与评分单元共用这一几何构造，确保评分约束与后验计算一致。

令读数步长为$h=360^\circ/N$，半步$\delta=h/2$。任意连续读数距最近格中心至多$\delta$，因此必须同步增宽所有随读数转动的约束：
\begin{equation}
 \begin{aligned}
 \text{扇区半角}&:\ \varepsilon+\delta,\qquad
 \text{前向投影下界}:\ 5\cos(\varepsilon+\delta),\\
 \text{条带半宽}&:\ R_2\sin\varepsilon+2R_{\rm out}\sin(\delta/2),\\
 R_2&=\min(1500,r_{2,\max}),\qquad
 R_{\rm out}=\max_{v\in\operatorname{vert}Q(q)}\|v-S(q)\|.
 \end{aligned}
 \label{eq:half_step}
\end{equation}
条带增量来自旋转前后单位法向量的差范数至多$2\sin(\delta/2)$。投影下界也须随角度放松：若原相对角为$t\in[-\varepsilon,\varepsilon]$且$r\cos t\ge5\cos\varepsilon$，读数旋转$|\Delta|\le\delta$后的投影至少为$5\cos(\varepsilon+|\Delta|)\ge5\cos(\varepsilon+\delta)$。该推导同时包含原始外包中因近区放松引入的点。以增宽约束形成第$k$个单元外包$P_k(q)$，便有$U(q,\psi)\subseteq P_k(q)$，进而
\begin{equation}
 J(q)\le B(q)=\max_k\operatorname{diam}P_k(q).
 \label{eq:geometric_bound}
\end{equation}
该上界直接由$U(q,\psi)\subseteq P_k(q)$及直径关于集合包含关系的单调性得到。
仅取未经增宽的有限读数样本最大值，不能作为连续最坏上界。默认最终读数步长为$0.1^\circ$，评分扇区半角为$1.05^\circ$；实际观测后的扇区仍使用$1^\circ$。

\subsubsection{评分组合与分层选点}
两种方法分别给出同一最坏直径目标的上界，故取二者较小值作为最终评分：
\begin{equation}
 C(q)=\min\{A(q),B(q)\}\ge J(q).
 \label{eq:combined_bound}
\end{equation}
较小的$C$表示较紧的最坏保证，未必意味着未知精确值$J$也按同样顺序排列。

在候选搜索层面，借鉴DIRECT算法\cite{jones1993}与多层坐标搜索（MCS）\cite{huyer1999}兼顾全域探索和局部细化的思想，本文采用极坐标粗网格搜索，并围绕当前优选点及其镜像方位逐级细化。具体实现采用预设步长和邻域扩充，未使用DIRECT的潜在最优矩形筛选或MCS的层级坐标分裂机制；因此仅借用搜索思想，不沿用原算法的全局收敛保证。

在式\eqref{eq:p2_feasible}内依次采用100米/$5^\circ$、20米/$1^\circ$及5米/$0.25^\circ$的移动距离与方位步长搜索，并加入各方位的安全边界点。各层保留既有候选，最终按统一读数精度复核。候选点按精度优先、距离次优的字典序筛选\cite{mavrotas2009}：先比较$C(q)$，仅评分平局时优先较短移动；实现按$10^{-7}$米舍入判定数值平局。记推荐点为$q_*$，给定可调容差$\eta$，近优候选区域定义为
\begin{equation}
 \mathcal C_{\rm good}=\{q\in\mathcal F:C(q)\le(1+\eta)C(q_*)\},\qquad \eta=10\%.
\end{equation}
数值输出给出该连续区域中已访问的代表点，未证明连续全局最优。
另以$\mathcal V$表示已访问候选集，提供
\begin{equation}
 q_\eta\in\arg\min_{q\in\mathcal V\cap\mathcal C_{\rm good}}\|q\|.
 \label{eq:p2_epsilon_selection}
\end{equation}
这采用$\epsilon$-约束思想\cite{mavrotas2009}，将定位要求写成阈值，再缩短路程；文献中的$\epsilon$指目标阈值，与测向误差$\varepsilon$不同。由于本次移动与检测时间为$\|q\|/5+5$，缩短路程等价于减少该项时间。默认$\eta=10\%$是可调整偏好，只允许评分上界放宽，并非实际误差增加比例或文献给定最优值；它不替代默认精度优先方案。当前程序未实现AUGMECON，不据此声称连续Pareto最优。

计算时先求$A$，依次扫描$B$的单元直径。已扫最大值$b_t$单调不减，一旦$b_t\ge A$，便已确定$C=A$，可提前停止；此时完整$B$与全局最坏单元仍未知，$b_t$只作为$B$的下界记录。不能因某点的$A$较大而淘汰它，因为$A$不是评分下界。

\begin{table}[htbp]
 \centering
 \caption{第二问自建场景的同口径比较（均保证收信）}
 \label{tab:p2_selection}
 \begin{tabular}{llrrr}
 \hline
 首点/米 & 策略 & 局部第二点/米 & $C$/米 & 移动加检测/秒\\
 \hline
 $(0,0)$ & 解析基线 & $(750,660)$ & 115.6930 & 204.810\\
 $(0,0)$ & 本文选点 & $(842.144,546.895)$ & 112.3293 & 205.828\\
 $(0,0)$ & 10\%内最短点 & $(773.511,541.618)$ & 122.8653 & 193.857\\
 $(1500,0)$ & 解析基线 & $(750,660)$ & 54.0236 & 204.810\\
 $(1500,0)$ & 本文选点 & $(220.942,195.473)$ & 15.7688 & 64.000\\
 \hline
 \end{tabular}
\end{table}
表\ref{tab:p2_selection}的首示向度均为$0^\circ$，圆外包取360边、读数步长取$0.1^\circ$，时间仅计本次移动及5秒检测。相对于解析基线点的同口径评分，原点和边界场景的上界分别下降2.91\%和70.81\%；这表示最坏直径保证变紧，不等于真实定位误差同比下降。原点场景在1281个自建源组合中全部收信；安全保证仍以四圆盘证明为依据。

默认原点流程访问548个候选，完整扫描与提前停止得到相同的逐层优选点、最终评分及排序。扫描量由3\,205\,440降至3\,201\,498，仅节省0.123\%；本次完整流程耗时分别为14.821秒与14.868秒，因此不能据此声称显著提速。固定本文原点选点，将精度提高至720边、$0.05^\circ$后，上界由112.3293降至111.7932米，反映外包与读数网格仍有保守性。所有结果均来自自建离线实验；几何裁剪使用双精度浮点运算，尚非严格区间算术认证。
